\documentclass[11pt,a4paper]{article}

% ============================================================
% Standard English mathematical typesetting preamble
% ============================================================
\usepackage[utf8]{inputenc}
\usepackage[T1]{fontenc}
\usepackage[english]{babel}
\usepackage{amsmath,amssymb,amsthm}
\usepackage{geometry}
\geometry{margin=1in}
\usepackage{graphicx}
\usepackage{booktabs}
\usepackage{array}
\usepackage{fancyhdr}
\usepackage{titlesec}
\usepackage{enumitem}
\usepackage{mathtools}
\usepackage{bm}

% Theorem-like environments
\theoremstyle{plain}
\newtheorem{theorem}{Theorem}
\newtheorem{lemma}{Lemma}
\theoremstyle{definition}
\newtheorem{remark}{Remark}

% Page style
\pagestyle{fancy}
\fancyhf{}
\fancyhead[C]{\small Cayley --- Collected Mathematical Papers, Vol.\ VIII}
\fancyfoot[C]{\thepage}
\renewcommand{\headrulewidth}{0.4pt}

% Section formatting to match original numbered articles
\titleformat{\section}{\large\scshape\centering}{}{0em}{}
\titleformat{\subsection}{\normalsize\bfseries}{}{0em}{}

% Custom commands for Cayley's notation
\newcommand{\dd}{\,d}
\newcommand{\sourcepage}[1]{%
  \par\noindent\textit{[Source: Collected Mathematical Papers, Vol. VIII, book p.~#1.]}\par\medskip
}
\DeclareMathOperator{\Jac}{Jac}

\begin{document}

% ============================================================
% Tail of last paper (continuation from p.~566) -- book p.~567
% ============================================================

\noindent
\textit{p.~567 [555] (tail of paper on a system of cubic equations)}
\smallskip

\noindent
$a^2 d - 3 abc + 2 b^3$; the general value of $u$ is: any function whatever of the
last-mentioned three functions.  We have to find the rational non-integral functions
of these functions which are rational and integral functions of the coefficients; such a
function is
\[
\frac{1}{a^2}\{(a^2 d - 3abc + 2b^3) + \tfrac{1}{4}(ac - b^2)^3\}
   = a c^3 + 4 a c^2 + 4 b^3 d - 3 b^2 c^2 - 6 a b c d,
\]
and the original three solutions, together with the last-mentioned function $a^2 d^2 + \&c.$,
constitute the complete system of the seminvariants of the cubic function; viz.\ every
other seminvariant is a rational \emph{and integral} function of these.  And so, in the general
case, the problem is to complete the series of the $n$ solutions $a$, $ac-b^2$, $a^2 d - 3abc + 2b^3$,
$a^3 e - 4a^2 bd + 6ab^2 c - 3b^4$, \&c. by adding thereto the solutions which, being rational but
non-integral functions of these, are rational and integral functions of the coefficients;
and thus to arrive at a series of solutions such that every other solution is a rational
and integral function of these.

\vspace{1em}
\hrule height 0.4pt width 0.3\linewidth
\vspace{1em}

\subsection*{4.\quad \textit{Note on certain Families of Surfaces.}\ \ Report, 1871, pp.~19, 20.}

See the paper numbered 526, of which this Note is a duplicate.

\vspace{1em}
\hrule height 0.4pt width 0.3\linewidth
\vspace{1em}

\subsection*{5.\quad \textit{On the Mercator's Projection of a Surface of Revolution.}\ \ Report, 1873, p.~9.}

\textsc{The} theory of Mercator's projection is obviously applicable to any surface of
revolution; the meridians and parallels are represented by two systems of parallel lines
at right angles to each other, in such wise that for the infinitesimal rectangles
included between two consecutive arcs of meridian and arcs of parallel the rectangle in
the projection is similar to that on the surface.  Or, what is the same thing, drawing
on the surface the meridians at equal infinitesimal intervals of angular distance, we
may draw the parallels at such intervals as to divide the surface into infinitesimal
squares; the meridians and parallels are then the projections represented by two
systems of equidistant parallel lines dividing the plane into squares.  And if the
angular distance between two consecutive meridians instead of being infinitesimal is
taken moderately small ($5^{\circ}$ or even $10^{\circ}$), then it is easy on the surface or \emph{in plano},
using only the curve which is the meridian of the surface, to lay down graphically
the series of parallels which are in the projection represented by equidistant parallel
lines.  The method is, of course, an approximate one, by reason that the angular distance
between the two consecutive meridians is finite instead of infinitesimal.

\bigskip

I have in this way constructed the projection of a skew hyperboloid of revolution;
viz.\ in one figure I show the hyperbola, which is the meridian section, and by means
of it (taking the interval of the meridians to be $= 10^{\circ}$) construct the positions of the

\clearpage

% ============================================================
% Book p.~568 -- continuation of Mercator's projection
% ============================================================

\noindent
\textit{p.~568 [555]}
\smallskip

\noindent
successive parallels; I complete the figure by drawing the hyperbolas which are the
orthographic projections of the meridians, and the right lines which are the ortho-
graphic projections of the parallels; the figure thus exhibits the orthographic projection
(on the plane of a meridian) of the hyperboloid divided into squares as above.  The
other figure, which is the Mercator's projection, is simply two systems of equidistant
parallel lines dividing the paper into squares.  I remark that in the first figure the
projections of the right lines on the surface are the tangents to the bounding hyper-
bola; in particular, the projection of one of these lines is an asymptote of the
hyperbola.  This I exhibit in the figure, and by means of it trace the Mercator's pro-
jection of the right line on the surface; viz.\ this is a serpentine curve included
between the right lines which represent two opposite meridians and having these lines
for asymptotes.  It is sufficient to draw one of these curves, since obviously for any
other line of the surface belonging to the same system the Mercator's projection is
at once obtained by merely displacing the curve parallel to itself; and for any line of
the other system the projection is a like curve in a reversed position.

\bigskip

A Mercator's projection might be made of a skew hyperboloid not of revolution;
viz.\ the curves of curvature ought be drawn so as to divide the surface into squares,
and the curves of curvature be then represented by equidistant parallel lines as above;
and the construction would in only a little more difficult.  The projection presented
itself to me as a convenient one for the representation of the geodesic lines on the
surface, and for exhibiting them in relation to the right lines of the surface; but I
have not yet worked this out.  In conclusion, it may be remarked that a surface in
general cannot be divided into squares by its curves of curvature, but that it may
be in an infinity of ways divided into squares by two systems of curves on the
surface, and any such system of curves gives rise to a Mercator's projection of the
surface.

\clearpage

% ============================================================
% Book pp. 569-570 -- source-checked replacement
% ============================================================

\sourcepage{569}

\begin{center}
  {\Large\scshape Notes and References.}
\end{center}

\noindent
518. \textsc{Ribaucour}, \textit{C.~R.}, t.~\textsc{LXXV}. (1872), pp.~533--536, referring to my Note remarks
that the condition can be (by means of the imaginary coordinates of M.~Ossian Bonnet)
expressed in a simple form communicated by him to the Philomathic Society, May,
1870. I reproduce this investigation, although it is not easy to present it in a quite
intelligible form. We take $p=f(x,y,z)$ to represent a family of surfaces belonging to
a triply orthotomic system, and consider two neighbouring surfaces $(A)$ and $(A')$ corresponding
to the values $z$ and $z+dz$; $A$ and $A'$ the two points where they meet the
trajectories of the surfaces; $AT$, $A'T'$ the tangents to the curves of curvature of the
same system at $A$, $A'$ respectively. Then according to the remark of M.~L\'evy, it is
to be expressed that these lines meet, and this is done by expressing that along the
trajectory $AA'$, the variation of the angle of $AT$ with the osculating plane at $A$ is
equal to the angle of the osculating planes at $A$, $A'$ respectively.

Let $B'$ be the spherical image of $A'$; the plane $OBB'$ is parallel to the osculating
plane at $A$ of the trajectory, and the angle of the two osculating planes measures
the geodesic curvature of $BB'$: denote this by $d\gamma$.

Let $\beta$ be the angle of $BB'$ with $BX$, $\theta$ the angle of $AT$ with $BX$, $\beta-\theta$ is the
angle of $AT$ with the osculating plane at $A$ of the trajectory: $d\beta-d\theta=d\gamma$. Introducing
the symmetric imaginary coordinates $x$ and $y$, we write
\[
  a=\frac{dp}{\lambda^2\,dx},\qquad
  b=\frac{dp}{\lambda^2\,dy},\qquad
  c=\frac{1}{\lambda^2}\frac{d^2p}{dx\,dy},\qquad
  ds^2=4\lambda^2\frac{da}{dx}\frac{db}{dy}\,dx\,dy.
\]

But $dx$ and $dy$ being the increments of $x$, $y$ corresponding to $dz$ in the passage
from $A$ to $A'$, then by a theorem of M.~Liouville
\[
  d\gamma=d\beta-i\left(\frac{d\lambda}{\lambda\,dx}\,dx
  -\frac{d\lambda}{\lambda\,dy}\,dy\right);
\]
the condition thus is
\[
  d\theta=i\left(\frac{d\lambda}{\lambda\,dx}\,dx
  -\frac{d\lambda}{\lambda\,dy}\,dy\right),
\]

\vfill
\noindent C. VIII.\hfill 72

\clearpage

\sourcepage{570}
\noindent
\hbox to \textwidth{570\hfill{\scshape Notes and References.}\hfill}

\medskip
\noindent
and the formula
\[
  e^{-2i\theta}=\pm\sqrt{\frac{da}{dx}}\div\sqrt{\frac{db}{dy}},
\]
enables this to be written in the definitive form
\[
\begin{aligned}
&dx\,\frac{d}{dx}\,i\left(\lambda^4\frac{db}{dy}\div\frac{da}{dx}\right)
 +dy\,\frac{d}{dy}\,i\left(\frac{db}{dy}\div\lambda^4\frac{da}{dx}\right)\\
&\hspace{7em}
 +dz\left\{\frac{d}{dz}\left(i\frac{db}{dy}\right)
       -\frac{d}{dz}\left(i\frac{da}{dx}\right)\right\}=0.
\end{aligned}
\]
We have
\[
  dx\left(\frac{1}{2}p+c\right)+dy\,\frac{db}{dy}+dz\,\frac{db}{dz}=0,
\]
\[
  dx\,\frac{da}{dx}+dy\left(\frac{1}{2}p+c\right)+dz\,\frac{da}{dz}=0,
\]
and thence eliminating $dx$, $dy$, $dz$, we have
\[
\left|
\begin{array}{ccc}
\dfrac{d}{dx}\,i\left(\lambda^4\dfrac{db}{dy}\div\dfrac{da}{dx}\right) &
\dfrac{d}{dy}\,i\left(\dfrac{db}{dy}\div\lambda^4\dfrac{da}{dx}\right) &
\dfrac{d}{dz}\,i\left(\dfrac{db}{dy}\div\dfrac{da}{dx}\right) \\[1.2em]
\dfrac{1}{2}p+c & \dfrac{db}{dy} & \dfrac{db}{dz} \\[1.2em]
\dfrac{da}{dx} & \dfrac{1}{2}p+c & \dfrac{da}{dz}
\end{array}
\right|=0,
\]
which defines the triply orthotomic system.

\vfill
\begin{center}
{\scshape End of Vol. VIII.}
\end{center}

\vfill
\begin{center}
\rule{0.48\linewidth}{0.3pt}\\[0.5em]
\footnotesize CAMBRIDGE: PRINTED BY J. AND C. F. CLAY, AT THE UNIVERSITY PRESS.
\end{center}

\end{document}

